\chapter{Literature Review}
\setcounter{equation}{0}
\section{GramBeddings}
Bozkir, Dalgi\c{c}, and Aydos propose GramBeddings, a neural-network method for identifying phishing web pages from URLs. The method represents URL text using character-level n-grams and learned embeddings, allowing the model to capture local patterns in strings such as unusual tokens, word fragments, and character sequences. The work demonstrates the value of neural representations for URL-only phishing detection on a large dataset.\cite{bib2}

The paper is relevant to Misiek and Hyla's base paper because GramBeddings is one of the datasets used in the X-PHIDE cross-dataset evaluation. It also provides a meaningful neural comparison for X-PHIDE: both methods use only URL-derived information, but X-PHIDE emphasises a small set of stable, engineered features and the efficiency of XGBoost, whereas GramBeddings learns representations directly from URL text.\cite{bib1,bib2}

\section{Unsupervised Domain Adaptation}
Rashid, Doyle, Han, and Seneviratne study the generalisation of phishing URL detection using unsupervised domain adaptation. Their work addresses the situation in which labelled training data are available in a source domain but the target domain has different characteristics and limited or no labels. The approach aims to reduce the effect of domain shift so that a detector can transfer more reliably between URL datasets.\cite{bib3}

This paper is directly relevant to the base paper because both focus on performance across changing data distributions rather than performance on a single random split. X-PHIDE does not primarily adapt a model with unlabelled target data; instead, it investigates cross-dataset performance, selects features that remain useful across datasets, and evaluates the resulting XGBoost model on independent sources. Domain adaptation therefore provides a complementary strategy and reinforces the base paper's motivation for testing generalisation explicitly.\cite{bib1,bib3}

\section{Survey of Phishing Environments and Countermeasures}
Aleroud and Zhou present a survey of phishing environments, attack techniques, and countermeasures. The paper organises the broader phishing problem by examining how attacks are delivered, how users are deceived, and how technical and organisational controls can reduce the threat. It provides background for understanding why detection systems must consider changing attacker behaviour and diverse sources of evidence.\cite{bib4}

The survey is relevant to the base paper as a broad security context for URL-based detection. While it covers phishing countermeasures more widely than X-PHIDE, its discussion of varied attack environments supports the base paper's concern that a detector trained under one collection condition may not generalise to another. X-PHIDE addresses this concern more narrowly through URL-only features, cross-dataset validation, and a browser plug-in workflow.\cite{bib1,bib4}

\section{Machine Learning and Deep Learning Review}
Divakaran and Oest review the use of machine learning and deep learning for phishing detection. The review compares families of detection methods, input features, learning architectures, and evaluation practices, showing how conventional machine-learning models and neural networks have both been applied to phishing-related classification tasks. It also highlights the importance of selecting appropriate data and evaluation settings when comparing reported results.\cite{bib5}

This review is relevant to the base paper because it places X-PHIDE among the larger set of machine-learning and deep-learning approaches to phishing detection. The base paper follows the review's broad comparison perspective but focuses on a specific unresolved issue: whether strong results transfer across independent URL datasets. Its comparison of XGBoost with a neural model, together with measurements of training time, inference time, and memory use, gives the general review criteria a concrete application.\cite{bib1,bib5}

\section{Hybrid DNN--LSTM Detection}
Özcan, Catal, Dönmez, and \c{S}entürk propose a hybrid DNN--LSTM model for detecting phishing URLs. The model combines dense neural-network processing with LSTM-based sequence modelling so that both feature relationships and sequential patterns in URL data can be learned. This represents a more complex neural alternative to traditional classifiers and illustrates how URL detection can use learned representations rather than relying only on manually selected features.\cite{bib6}

The hybrid model is relevant to the base paper because Misiek and Hyla compare X-PHIDE with a neural-network approach while considering practical deployment costs. The two studies therefore address a similar detection task with different modelling priorities: the DNN--LSTM work emphasises neural sequence modelling, while X-PHIDE emphasises stable cross-dataset features, XGBoost classification, and browser-level efficiency. This comparison helps justify evaluating accuracy together with computational requirements.\cite{bib1,bib6}
